\begin{proofbox}
    \textbf{\textcolor{CProof}{统一题设与证明目标}}

    已知抛物线 $x^2=2py\;(p>0)$，$A(x_1,y_1)$、$B(x_2,y_2)$ 为曲线上两个不同点，
    过 $A,B$ 的切线交于 $P$；$M$ 为 $AB$ 中点。
    需要证明
    \[
        P\left(\frac{x_1+x_2}{2},\frac{x_1x_2}{2p}\right),\qquad
        PM\parallel y\text{轴},\qquad
        S_{\triangle PAB}=\frac{|x_1-x_2|^3}{8p}.
    \]
    以下四种方法均只使用高中范围内的直线方程、二次函数、平方差公式、
    中点坐标、三角形面积等知识；最后补充切点弦及附加几何性质。

    \medskip

    \textbf{\textcolor{CProof}{证明前的初中铺垫：先知道为什么要找中线}}

    \textbf{普通三角形：}若 $M$ 是 $AB$ 中点，且 $P,M$ 在同一条竖直线上，
    则 $AM=MB$，$PM$ 将整个三角形分成两块。
    以 $PM$ 为共同底边，两块三角形的高分别为
    $|x_A-x_P|$、$|x_B-x_P|$（注意：这是\textbf{水平距离}）。因此
    \[
        S_{\triangle PAB}
        =\frac12 PM\,|x_A-x_P|+\frac12 PM\,|x_B-x_P|.
    \]
    又因为 $x_P=x_M=(x_A+x_B)/2$，所以
    \[
        |x_A-x_P|=|x_B-x_P|=\frac{|x_A-x_B|}{2},
        \quad\Longrightarrow\quad
        \boxed{S_{\triangle PAB}=\frac12 PM\,|x_A-x_B|}.
    \]
    \textbf{圆的切线三角形：}由 $PA=PB$ 可得 $PM\perp AB$，
    所以圆可以直接以 $AB$ 为底、$PM$ 为高。
    \textbf{抛物线的切线三角形：}一般不是等腰三角形，
    这里先要用代数\textbf{证明 $PM$ 竖直}，然后才可使用上面的分割面积方法。
    这样，“为什么相减求 $x_P$”“为什么求中点 $M$”就有了清楚的目标。

    \medskip

    \par\noindent\begin{minipage}{\linewidth}
    \textbf{\textcolor{CProof}{方法一：门槛最低——竖直中线分成两个三角形}}

    \begin{center}
\begin{tikzpicture}[x=1.05cm,y=1.05cm,font=\small,>=Stealth]
\path (-2,-2.1) rectangle (9,3.6);
\coordinate (A) at (-1,.25); \coordinate (B) at (3,2.25);
\coordinate (P) at (1,-.75); \coordinate (M) at (1,1.25);
\coordinate (N) at (1,.25); \coordinate (H) at (.2,.85);
\fill[blue!12] (P)--(A)--(M)--cycle;
\fill[orange!20] (P)--(M)--(B)--cycle;
\draw[->,gray] (-1.8,0)--(3.8,0) node[right] {$x$};
\draw[->,gray] (0,-1.35)--(0,3.2) node[above] {$y$};
\draw[blue!70!black,thick,domain=-1.6:3.45,samples=100] plot (\x,{\x*\x/4});
\draw[gray!80,domain=-1.7:1.5] plot (\x,{-.5*\x-.25});
\draw[gray!80,domain=.7:3.4] plot (\x,{1.5*\x-2.25});
\draw[very thick] (A)--(B);
\draw[thick] (A)--(P)--(B);
\draw[orange!80!black,very thick] (P)--(M);
\foreach \pt/\pos in {B/above left,M/above left}
  \fill (\pt) circle (1.6pt) node[\pos] {$\pt$};
\fill (A) circle (1.6pt); \node[anchor=east] at (-1.3,-.35) {$A$};
\fill (P) circle (1.6pt); \node[anchor=west] at (1.4,-.65) {$P$};

\draw[orange!80!black,dashed] (1,-.9)--(1,2.7);
\draw[teal!80!black] (-1,.25)--(1,.25);
\draw[dashed,teal!80!black] (-1,.25)--(-1,-1.5) (1,-.9)--(1,-1.5);
\draw[<->,teal!80!black] (-1,-1.45)--(1,-1.45) node[midway,below] {$h_A$};
\draw[<->,teal!80!black] (1,2.25)--(3,2.25) node[midway,above] {$h_B$};
\draw (1,.4)--(.85,.4)--(.85,.25);
\draw (1,2.1)--(1.15,2.1)--(1.15,2.25);

\node[anchor=north west,align=left,text width=4.7cm] at (4.1,3.1) {\textbf{同底 $PM$，水平作高}\\[8pt]
$h_A=|x_1-x_P|$\\[5pt]
$h_B=|x_2-x_P|$\\[8pt]
$h_A+h_B=|x_2-x_1|$\\[8pt]
$S_{PAM}+S_{PMB}$\\[4pt]
$\quad=\dfrac12 PM(h_A+h_B)$\\[8pt]
水平高落在直线 $PM$ 上，\\不必落在线段 $PM$ 内。};
\end{tikzpicture}
\end{center}

    \end{minipage}\par

    \textbf{思路：}不用点到斜线的距离公式。先求 $P$，再求中点 $M$；
    把 $PM$ 作为\textbf{竖直的底边}，左右两块三角形分别按“底乘高除以二”计算。

    \begin{description}[style=nextline, leftmargin=2.8em, labelsep=0.5em, itemsep=0.55em]
        \item[\textbf{第一步（写出两条切线）：}]
        因为 $A,B$ 在抛物线上，
        \[
            y_1=\frac{x_1^2}{2p},\qquad y_2=\frac{x_2^2}{2p}.
        \]
        使用抛物线在 $(x_i,y_i)$ 处的切线点式 $x_ix=p(y+y_i)$，得
        \[
            \ell_A:\ x_1x=p(y+y_1),\qquad
            \ell_B:\ x_2x=p(y+y_2).
        \]
        \textbf{为什么它是切线？}以 $A$ 为例，直线可改写为
        \[
            y=\frac{x_1}{p}x-\frac{x_1^2}{2p}.
        \]
        同一横坐标处，抛物线的纵坐标减去直线的纵坐标为
        \[
            \frac{x^2}{2p}-\left(\frac{x_1}{p}x-\frac{x_1^2}{2p}\right)
            =\frac{(x-x_1)^2}{2p}\ge 0.
        \]
        只有 $x=x_1$ 才等号成立，说明直线恰在 $A$ 处与抛物线接触。
        对 $B$ 同理。这也提供了一个\textbf{不需要导数}的切线检验方法。

        \item[\textbf{第二步（两式相减，求交点横坐标）：}]
        设 $P=(x_P,y_P)$。因为 $P$ 在两条切线上，
        \[
            x_1x_P=p(y_P+y_1),\qquad
            x_2x_P=p(y_P+y_2).
        \]
        \textbf{为什么相减？}两式中都有 $py_P$，相减可把未知的纵坐标消去：
        \[
        \begin{aligned}
            (x_1-x_2)x_P
            &=p(y_1-y_2)\\
            &=p\left(\frac{x_1^2}{2p}-\frac{x_2^2}{2p}\right)\\
            &=\frac{x_1^2-x_2^2}{2}\\
            &=\frac{(x_1-x_2)(x_1+x_2)}2.
        \end{aligned}
        \]
        由于 $A\ne B$，故 $x_1\ne x_2$，可以除以 $x_1-x_2$，得到
        \[
            \boxed{x_P=\frac{x_1+x_2}{2}}.
        \]
        \textbf{几何意义：}交点 $P$ 的横坐标是两切点横坐标的平均数。

        \item[\textbf{第三步（回代，求交点纵坐标）：}]
        将 $x_P=(x_1+x_2)/2$ 代入 $x_1x_P=p(y_P+y_1)$：
        \[
        \begin{aligned}
            py_P
            &=x_1\frac{x_1+x_2}{2}-py_1\\
            &=\frac{x_1^2+x_1x_2}{2}-\frac{x_1^2}{2}\\
            &=\frac{x_1x_2}{2},
        \end{aligned}
        \]
        因 $p>0$，可除以 $p$，于是
        \[
            \boxed{P\left(\frac{x_1+x_2}{2},\frac{x_1x_2}{2p}\right)}.
        \]

        \item[\textbf{第四步（找中点与竖直中线）：}]
        中点公式给出
        \[
        \begin{aligned}
            x_M&=\frac{x_1+x_2}{2}=x_P,\\
            y_M&=\frac{y_1+y_2}{2}
            =\frac12\left(\frac{x_1^2}{2p}+\frac{x_2^2}{2p}\right)
            =\frac{x_1^2+x_2^2}{4p}.
        \end{aligned}
        \]
        $P$ 和 $M$ 的横坐标相等，所以 $PM$ 竖直，平行于 $y$ 轴。
        \[
        \begin{aligned}
            y_M-y_P
            &=\frac{x_1^2+x_2^2}{4p}-\frac{x_1x_2}{2p}\\
            &=\frac{x_1^2-2x_1x_2+x_2^2}{4p}\\
            &=\frac{(x_1-x_2)^2}{4p}>0.
        \end{aligned}
        \]
        因此 $M$ 在 $P$ 上方，且
        \[
            \boxed{PM=\frac{(x_1-x_2)^2}{4p}}.
        \]

        \item[\textbf{第五步（把三角形分成左右两块求面积）：}]
        以竖直线段 $PM$ 为底，点 $A$ 到这条\textbf{竖直直线}的高为水平距离
        \[
            |x_1-x_P|=\left|x_1-\frac{x_1+x_2}{2}\right|
            =\frac{|x_1-x_2|}{2}.
        \]
        同理，点 $B$ 到直线 $PM$ 的高也是 $|x_1-x_2|/2$。
        由于 $M$ 在 $AB$ 上，$\triangle PAM$ 与 $\triangle PMB$ 恰好拼成 $\triangle PAB$，故
        \[
        \begin{aligned}
            S_{\triangle PAB}
            &=\frac12 PM\cdot|x_1-x_P|
             +\frac12 PM\cdot|x_2-x_P|\\
            &=\frac12\cdot\frac{(x_1-x_2)^2}{4p}
                \left(\frac{|x_1-x_2|}{2}+\frac{|x_1-x_2|}{2}\right)\\
            &=\boxed{\frac{|x_1-x_2|^3}{8p}}.
        \end{aligned}
        \]
        \textbf{几何意义：}“水平宽度”贡献一次方，“竖直中线”贡献平方，
        相乘便是三次方。\textbf{$PM$ 并不是一般意义下 $AB$ 边上的高。}
    \end{description}

    \medskip
    \par\noindent\begin{minipage}{\linewidth}
    \textbf{\textcolor{CProof}{方法二：速度最快——切线相减 $\to$ 中线 $\to$ 面积}}

    \begin{center}
\begin{tikzpicture}[x=1.05cm,y=1.05cm,font=\small,>=Stealth]
\path (-2,-2.1) rectangle (9,3.6);
\coordinate (A) at (-1,.25); \coordinate (B) at (3,2.25);
\coordinate (P) at (1,-.75); \coordinate (M) at (1,1.25);
\coordinate (N) at (1,.25); \coordinate (H) at (.2,.85);
\fill[blue!12] (P)--(A)--(M)--cycle;
\fill[orange!20] (P)--(M)--(B)--cycle;
\draw[->,gray] (-1.8,0)--(3.8,0) node[right] {$x$};
\draw[->,gray] (0,-1.35)--(0,3.2) node[above] {$y$};
\draw[blue!70!black,thick,domain=-1.6:3.45,samples=100] plot (\x,{\x*\x/4});
\draw[gray!80,domain=-1.7:1.5] plot (\x,{-.5*\x-.25});
\draw[gray!80,domain=.7:3.4] plot (\x,{1.5*\x-2.25});
\draw[very thick] (A)--(B);
\draw[thick] (A)--(P)--(B);
\draw[orange!80!black,very thick] (P)--(M);
\foreach \pt/\pos in {B/above left,M/above left}
  \fill (\pt) circle (1.6pt) node[\pos] {$\pt$};
\fill (A) circle (1.6pt); \node[anchor=east] at (-1.3,-.35) {$A$};
\fill (P) circle (1.6pt); \node[anchor=west] at (1.8,-1.15) {$P$};

\draw[dashed,teal!80!black] (M)--(3.65,1.25) (P)--(3.65,-.75);
\draw[<->,teal!80!black] (3.55,-.75)--(3.55,1.25);

\node[anchor=north west,align=left,text width=4.7cm] at (4.1,3.1) {\textbf{相减定位，再作纵坐标差}\\[8pt]
$x_P=x_M=\dfrac{x_1+x_2}{2}$\\[10pt]
$y_M=\dfrac{x_1^2+x_2^2}{4p}$\\[8pt]
$y_P=\dfrac{x_1x_2}{2p}$\\[10pt]
右侧箭头长度为 $y_M-y_P$，\\即竖直中线 $PM$。};
\end{tikzpicture}
\end{center}

    \end{minipage}\par

    \textbf{使用时机：}题目给出了 $x_1,x_2$，并且可以直接使用已证明的切线点式。
    保留考场上最有效的运算，不再重复切线公式的证明。

    \begin{description}[style=nextline, leftmargin=2.8em, labelsep=0.5em, itemsep=0.5em]
        \item[\textbf{第一步（相减确定 $P$）：}]
        交点满足 $x_1x_P=p(y_P+y_1)$、$x_2x_P=p(y_P+y_2)$。于是
        \[
            (x_1-x_2)x_P=p(y_1-y_2)
            =\frac{x_1^2-x_2^2}{2},
        \]
        即 $x_P=(x_1+x_2)/2$。代回第一式，得
        \[
            y_P=\frac{x_1x_P}{p}-\frac{x_1^2}{2p}
            =\frac{x_1x_2}{2p}.
        \]
        \item[\textbf{第二步（平方差求中线）：}]
        \[
            y_M=\frac{x_1^2+x_2^2}{4p},\qquad
            PM=y_M-y_P=\frac{(x_1-x_2)^2}{4p}.
        \]
        \textbf{原因：}$x_M=x_P$，所以中线竖直；$PM$ 的长度就是两纵坐标之差。
        \item[\textbf{第三步（代一次面积公式）：}]
        两个小三角形都以 $PM$ 为底，水平高之和为 $|x_1-x_2|$：
        \[
            \boxed{S=\frac12 PM\cdot|x_1-x_2|
            =\frac12\cdot\frac{(x_1-x_2)^2}{4p}|x_1-x_2|
            =\frac{|x_1-x_2|^3}{8p}}.
        \]
    \end{description}
    \textbf{考场识别信号：}看见“两切线交点 + 求面积”，优先用此法；
    省掉求 $AB$ 长度和点到 $AB$ 距离的根式运算。

    \medskip
    \par\noindent\begin{minipage}{\linewidth}
    \textbf{\textcolor{CProof}{方法三：步骤最少——以横坐标中点作对称代换}}

    \begin{center}
\begin{tikzpicture}[x=1.05cm,y=1.05cm,font=\small,>=Stealth]
\path (-2,-2.1) rectangle (9,3.6);
\coordinate (A) at (-1,.25); \coordinate (B) at (3,2.25);
\coordinate (P) at (1,-.75); \coordinate (M) at (1,1.25);
\coordinate (N) at (1,.25); \coordinate (H) at (.2,.85);
\fill[blue!12] (P)--(A)--(M)--cycle;
\fill[orange!20] (P)--(M)--(B)--cycle;
\draw[->,gray] (-1.8,0)--(3.8,0) node[right] {$x$};
\draw[->,gray] (0,-1.35)--(0,3.2) node[above] {$y$};
\draw[blue!70!black,thick,domain=-1.6:3.45,samples=100] plot (\x,{\x*\x/4});
\draw[gray!80,domain=-1.7:1.5] plot (\x,{-.5*\x-.25});
\draw[gray!80,domain=.7:3.4] plot (\x,{1.5*\x-2.25});
\draw[very thick] (A)--(B);
\draw[thick] (A)--(P)--(B);
\draw[orange!80!black,very thick] (P)--(M);
\foreach \pt/\pos in {B/above left,M/above left}
  \fill (\pt) circle (1.6pt) node[\pos] {$\pt$};
\fill (A) circle (1.6pt); \node[anchor=east] at (-1.3,-.35) {$A$};
\fill (P) circle (1.6pt); \node[anchor=west] at (1.4,-.65) {$P$};

\draw[dashed] (A)--(-1,-1.65) (B)--(3,-1.65) (P)--(1,-1.65);
\draw[<->] (-1,-1.55)--(1,-1.55) node[midway,below] {$t$};
\draw[<->] (1,-1.55)--(3,-1.55) node[midway,below] {$t$};



\node[anchor=north west,align=left,text width=4.7cm] at (4.1,3.1) {\textbf{对称的是横坐标}\\[8pt]
$x_P=x_M=x_0,\quad t>0$\\[8pt]
$y_M=\dfrac{x_0^2+t^2}{2p}$\\[8pt]
$y_P=\dfrac{x_0^2-t^2}{2p}$\\[8pt]
$PM=\dfrac{t^2}{p},\quad S=\dfrac{t^3}{p}$\\[8pt]
$A,B$ 的高度通常不同，\\并非关于 $x=x_0$ 镜像对称。};
\end{tikzpicture}
\end{center}

    \end{minipage}\par

    \textbf{思路：}不妨交换 $A,B$ 的名称使 $x_1<x_2$。令
    \[
        x_1=x_0-t,\qquad x_2=x_0+t,\qquad t>0.
    \]
    于是 $x_0=(x_1+x_2)/2$、$2t=x_2-x_1$。
    \textbf{几何意义：}$x_0$ 是底边中点的横坐标，$t$ 是两个切点水平距离的一半。

    \begin{description}[style=nextline, leftmargin=2.8em, labelsep=0.5em, itemsep=0.5em]
        \item[\textbf{第一步（用对称形式一次求出顶点与中点）：}]
        两条切线为
        \[
        \begin{cases}
            y=\dfrac{x_0-t}{p}x-\dfrac{(x_0-t)^2}{2p},\\[3pt]
            y=\dfrac{x_0+t}{p}x-\dfrac{(x_0+t)^2}{2p}.
        \end{cases}
        \]
        相减可得
        \[
            0=-\frac{2t}{p}x+\frac{2x_0t}{p},
            \quad t>0\ \Longrightarrow\ x=x_0.
        \]
        把 $x=x_0$ 代入任一切线：
        \[
        \begin{aligned}
            y_P
            &=\frac{(x_0-t)x_0}{p}-\frac{(x_0-t)^2}{2p}\\
            &=\frac{2x_0(x_0-t)-(x_0-t)^2}{2p}\\
            &=\frac{x_0^2-t^2}{2p}.
        \end{aligned}
        \]
        另由中点公式
        \[
            y_M=\frac{(x_0-t)^2+(x_0+t)^2}{4p}
            =\frac{x_0^2+t^2}{2p}.
        \]
        因而
        \[
            P\left(x_0,\frac{x_0^2-t^2}{2p}\right),\qquad
            M\left(x_0,\frac{x_0^2+t^2}{2p}\right).
        \]
        对照 $x_0=(x_1+x_2)/2$ 和
        $x_0^2-t^2=(x_0-t)(x_0+t)=x_1x_2$，同时得到顶点坐标公式。

        \item[\textbf{第二步（两个量相乘，直接出面积）：}]
        \[
            PM=y_M-y_P=\frac{t^2}{p},\qquad
            |x_2-x_1|=2t.
        \]
        以竖直中线为底、以两侧水平距离之和为高，得
        \[
            \boxed{S=\frac12\cdot\frac{t^2}{p}\cdot2t
            =\frac{t^3}{p}=\frac{|x_1-x_2|^3}{8p}}.
        \]
    \end{description}
    \textbf{方法特点：}只有“求 $P,M$”和“算面积”两个主要步骤；
    对称代换令平方和与平方差自然出现，但写出代换需要一点熟练度。

    \medskip
    \par\noindent\begin{minipage}{\linewidth}
    \textbf{\textcolor{CProof}{方法四：常规通用——切线方程、弦长、点线距离}}

    \begin{center}
\begin{tikzpicture}[x=1.05cm,y=1.05cm,font=\small,>=Stealth]
\path (-2,-2.1) rectangle (9,3.6);
\coordinate (A) at (-1,.25); \coordinate (B) at (3,2.25);
\coordinate (P) at (1,-.75); \coordinate (M) at (1,1.25);
\coordinate (N) at (1,.25); \coordinate (H) at (.2,.85);
\fill[blue!12] (P)--(A)--(M)--cycle;
\fill[orange!20] (P)--(M)--(B)--cycle;
\draw[->,gray] (-1.8,0)--(3.8,0) node[right] {$x$};
\draw[->,gray] (0,-1.35)--(0,3.2) node[above] {$y$};
\draw[blue!70!black,thick,domain=-1.6:3.45,samples=100] plot (\x,{\x*\x/4});
\draw[gray!80,domain=-1.7:1.5] plot (\x,{-.5*\x-.25});
\draw[gray!80,domain=.7:3.4] plot (\x,{1.5*\x-2.25});
\draw[very thick] (A)--(B);
\draw[thick] (A)--(P)--(B);
\draw[orange!80!black,very thick] (P)--(M);
\foreach \pt/\pos in {B/above left,M/above left}
  \fill (\pt) circle (1.6pt) node[\pos] {$\pt$};
\fill (A) circle (1.6pt); \node[anchor=east] at (-1.3,-.35) {$A$};
\fill (P) circle (1.6pt); \node[anchor=west] at (1.4,-.65) {$P$};

\draw[teal!80!black,very thick] (P)--(H);
\fill[teal!80!black] (H) circle (1.6pt) node[above] {$H$};
\draw[teal!80!black] ($(H)+(.161,.0805)$)--($(H)+(.2415,-.0805)$)--($(H)+(.0805,-.161)$);

\node[anchor=north west,align=left,text width=4.7cm] at (4.1,3.1) {\textbf{以 $AB$ 为底，须配 $PH$}\\[8pt]
$PH\perp AB,\quad d=PH$\\[8pt]
$S=\dfrac12 |AB|\,d$\\[10pt]
橙色 $PM$ 经过中点，\\绿色 $PH$ 垂直于底边。\\[8pt]
本图中 $H\ne M$，\\不能以 $PM$ 代替 $d$。};
\end{tikzpicture}
\end{center}

    \end{minipage}\par

    \textbf{思路：}这一路线适用于许多解析几何面积题：
    \textbf{求交点 $\to$ 求底边直线与长度 $\to$ 求顶点到该直线的距离 $\to$ 底乘高除以二。}
    它计算较长，但每一步都可迁移到其他题型。

    \begin{description}[style=nextline, leftmargin=2.8em, labelsep=0.5em, itemsep=0.55em]
        \item[\textbf{第一步（从相切条件推导切线方程）：}]
        设某条非竖直直线为 $y=kx+b$。将其代入抛物线方程，得到
        \[
            x^2=2p(kx+b)
            \ \Longleftrightarrow\ x^2-2pkx-2pb=0.
        \]
        切线与抛物线恰有一个公共点，所以此二次方程只有一个实数解，
        判别式为零：
        \[
            (-2pk)^2-4\cdot1\cdot(-2pb)
            =4p^2k^2+8pb=0.
        \]
        解得 $b=-pk^2/2$，且唯一解为
        \[
            x=\frac{2pk}{2}=pk.
        \]
        在 $A$ 处相切时 $x_1=pk$，即 $k=x_1/p$，于是
        \[
            y=\frac{x_1}{p}x-\frac{x_1^2}{2p},
            \quad\text{也就是}\quad x_1x=p(y+y_1).
        \]
        在 $B$ 处同理可得 $x_2x=p(y+y_2)$。
        \textbf{几何意义：}判别式等于零对应“直线只在一个位置碰到抛物线”。

        \item[\textbf{第二步（联立求两切线交点）：}]
        \[
            \begin{cases}
                x_1x_P=p(y_P+y_1),\\
                x_2x_P=p(y_P+y_2).
            \end{cases}
        \]
        两式相减：
        \[
            x_P=\frac{p(y_1-y_2)}{x_1-x_2}
            =\frac{x_1^2-x_2^2}{2(x_1-x_2)}
            =\frac{x_1+x_2}{2}.
        \]
        回代：
        \[
            y_P=\frac{x_1x_P}{p}-y_1
            =\frac{x_1(x_1+x_2)-x_1^2}{2p}
            =\frac{x_1x_2}{2p}.
        \]
        故 $P\left((x_1+x_2)/2,\ x_1x_2/(2p)\right)$。

        \item[\textbf{第三步（两点确定底边 $AB$ 的方程）：}]
        因 $x_2\ne x_1$，弦 $AB$ 的斜率为
        \[
        \begin{aligned}
            k_{AB}
            &=\frac{y_2-y_1}{x_2-x_1}
              =\frac{x_2^2-x_1^2}{2p(x_2-x_1)}\\
            &=\frac{x_1+x_2}{2p}.
        \end{aligned}
        \]
        过 $A$ 写点斜式：
        \[
            y-\frac{x_1^2}{2p}
            =\frac{x_1+x_2}{2p}(x-x_1).
        \]
        两边乘以 $2p$，并移项化简：
        \[
            2py-x_1^2=(x_1+x_2)x-x_1(x_1+x_2),
        \]
        \[
            \boxed{(x_1+x_2)x-2py-x_1x_2=0}.
        \]

        \item[\textbf{第四步（求底边长度）：}]
        两点间距离公式给出
        \[
        \begin{aligned}
            |AB|
            &=\sqrt{(x_2-x_1)^2+(y_2-y_1)^2}\\
            &=|x_2-x_1|\sqrt{1+
              \left(\frac{y_2-y_1}{x_2-x_1}\right)^2}\\
            &=|x_2-x_1|\sqrt{1+
              \frac{(x_1+x_2)^2}{4p^2}}\\
            &=\boxed{\frac{|x_2-x_1|}{2p}
              \sqrt{(x_1+x_2)^2+4p^2}}.
        \end{aligned}
        \]
        最后一行用到了 $p>0$，因此 $\sqrt{4p^2}=2p$。

        \item[\textbf{第五步（求顶点到 $AB$ 的垂直距离）：}]
        对直线 $ax+by+c=0$，点 $(u,v)$ 到该直线的距离为
        $|au+bv+c|/\sqrt{a^2+b^2}$。代入 $P$ 与 $AB$：
        \[
            d=\frac{\left|(x_1+x_2)\dfrac{x_1+x_2}{2}
              -2p\dfrac{x_1x_2}{2p}-x_1x_2\right|}
              {\sqrt{(x_1+x_2)^2+4p^2}}.
        \]
        分子逐项化简：
        \[
        \begin{aligned}
            \left|\frac{(x_1+x_2)^2}{2}-2x_1x_2\right|
            &=\left|\frac{x_1^2+2x_1x_2+x_2^2-4x_1x_2}{2}\right|\\
            &=\frac{(x_1-x_2)^2}{2}.
        \end{aligned}
        \]
        因为平方非负，取绝对值后数值不变，所以
        \[
            \boxed{d=\frac{(x_1-x_2)^2}
              {2\sqrt{(x_1+x_2)^2+4p^2}}}.
        \]
        \textbf{几何意义：}这里的 $d$ 才是以斜边 $AB$ 为底时真正的高，
        与前面作为竖直中线的 $PM$ 通常不同。

        \item[\textbf{第六步（底乘高求面积并验证中线）：}]
        \[
        \begin{aligned}
            S_{\triangle PAB}
            &=\frac12\cdot|AB|\cdot d\\
            &=\frac12\cdot\frac{|x_2-x_1|}{2p}
               \sqrt{(x_1+x_2)^2+4p^2}\cdot
               \frac{(x_1-x_2)^2}{2\sqrt{(x_1+x_2)^2+4p^2}}\\
            &=\boxed{\frac{|x_1-x_2|^3}{8p}}.
        \end{aligned}
        \]
        中点 $M$ 的横坐标是 $(x_1+x_2)/2=x_P$，故 $PM$ 平行 $y$ 轴。
    \end{description}

    \medskip
    \par\noindent\begin{minipage}{\linewidth}
    \textbf{\textcolor{CProof}{补充一：切点弦方程与切线存在条件}}

    \begin{center}
\begin{tikzpicture}[x=.72cm,y=.72cm,font=\small,>=Stealth]
\foreach \sh/\v/\ttl/\desc in {0/-.75/{两个切点}/{x_0^2-2py_0>0},5/.25/{一个切点}/{x_0^2-2py_0=0},10/1.25/{没有实切点}/{x_0^2-2py_0<0}} {
\begin{scope}[xshift=\sh cm]
\node at (1,3.3) {\ttl};
\draw[->,gray] (-1.8,0)--(3.6,0) node[right] {$x$};
\draw[->,gray] (0,-1.4)--(0,2.8);
\draw[blue!70!black,thick,domain=-1.6:3.2,samples=80] plot (\x,{\x*\x/4});
\ifdim\v pt<0pt
\draw[thick] (-1,.25)--(1,-.75)--(3,2.25);
\fill (-1,.25) circle (1.6pt); \node[anchor=east] at (-1.3,-.35) {$A$};
\fill (3,2.25) circle (1.6pt) node[above left] {$B$};
\node[orange!80!black,anchor=west] at (1.4,-.65) {$P$};
\else\ifdim\v pt<1pt
\draw[thick,domain=-1:3.1] plot (\x,{.5*\x-.25});
\node[orange!80!black,above] at (1,.65) {$P$};
\else
\node[orange!80!black,right] at (1,1.25) {$P$};
\fi\fi
\fill[orange!80!black] (1,\v) circle (2pt);
\node at (1,-1.9) {$\desc$};
\end{scope}}
\end{tikzpicture}
\end{center}

    \end{minipage}\par

    设 $P=(x_0,y_0)$。因为 $P$ 在 $A$ 处的切线上，
    \[
        x_0x_1=p(y_0+y_1)
        =py_0+\frac{x_1^2}{2}.
    \]
    移项得 $x_1^2-2x_0x_1+2py_0=0$，同理
    $x_2^2-2x_0x_2+2py_0=0$。因此 $x_1,x_2$ 是
    \[
        \boxed{x^2-2x_0x+2py_0=0}
    \]
    的两个不同实数解，根与系数的关系给出
    \[
        x_1+x_2=2x_0,\qquad x_1x_2=2py_0.
    \]
    将 $AB$ 的方程
    $(x_1+x_2)x-2py-x_1x_2=0$ 中的和与积替换，得
    \[
        2x_0x-2py-2py_0=0
        \ \Longrightarrow\ \boxed{x_0x=p(y+y_0)}.
    \]
    二次方程有两个不同实数解，当且仅当
    \[
        \Delta=(-2x_0)^2-4(2py_0)
        =4(x_0^2-2py_0)>0,
    \]
    即 $x_0^2-2py_0>0$。此时
    \[
        x_{1,2}=x_0\pm\sqrt{x_0^2-2py_0},
        \qquad |x_1-x_2|=2\sqrt{x_0^2-2py_0}.
    \]
    于是还能得到只含交点坐标的面积式
    \[
        \boxed{S_{\triangle PAB}
        =\frac{(x_0^2-2py_0)\sqrt{x_0^2-2py_0}}{p}}.
    \]
    \textbf{几何意义：}$x_0^2-2py_0>0$ 与
    $y_0<x_0^2/(2p)$ 等价，意味着 $P$ 位于抛物线的\textbf{下方}。
    $P$ 恰在曲线上时只有一个切点，两条切线三角形不再成立；
    $P$ 在曲线上方时不存在所需的两个实切点。

    \medskip
    \par\noindent\begin{minipage}{\linewidth}
    \textbf{\textcolor{CProof}{补充二：竖直中线被抛物线交点平分}}

    \begin{center}
\begin{tikzpicture}[x=1.05cm,y=1.05cm,font=\small,>=Stealth]
\path (-2,-2.1) rectangle (9,3.6);
\coordinate (A) at (-1,.25); \coordinate (B) at (3,2.25);
\coordinate (P) at (1,-.75); \coordinate (M) at (1,1.25);
\coordinate (N) at (1,.25); \coordinate (H) at (.2,.85);
\fill[blue!12] (P)--(A)--(M)--cycle;
\fill[orange!20] (P)--(M)--(B)--cycle;
\draw[->,gray] (-1.8,0)--(3.8,0) node[right] {$x$};
\draw[->,gray] (0,-1.35)--(0,3.2) node[above] {$y$};
\draw[blue!70!black,thick,domain=-1.6:3.45,samples=100] plot (\x,{\x*\x/4});
\draw[gray!80,domain=-1.7:1.5] plot (\x,{-.5*\x-.25});
\draw[gray!80,domain=.7:3.4] plot (\x,{1.5*\x-2.25});
\draw[very thick] (A)--(B);
\draw[thick] (A)--(P)--(B);
\draw[orange!80!black,very thick] (P)--(M);
\foreach \pt/\pos in {B/above left,M/above left}
  \fill (\pt) circle (1.6pt) node[\pos] {$\pt$};
\fill (A) circle (1.6pt); \node[anchor=east] at (-1.3,-.35) {$A$};
\fill (P) circle (1.6pt); \node[anchor=west] at (1.4,-.65) {$P$};

\fill[blue!70!black] (N) circle (1.7pt) node[right] {$N$};
\draw[thick] (.88,-.25)--(1.12,-.25) (.88,.75)--(1.12,.75);

\node[anchor=north west,align=left,text width=4.7cm] at (4.1,3.1) {\textbf{平分的是线段 $PM$}\\[8pt]
$x_N=x_P=x_M$\\[8pt]
$y_N=\dfrac{(x_1+x_2)^2}{8p}$\\[10pt]
$PN=NM=\dfrac{PM}{2}$\\[8pt]
$P$ 是切线交点；\\$N$ 才是同横坐标的曲线点。\\[8pt]
两段曲线弧不据此等长。};
\end{tikzpicture}
\end{center}

    \end{minipage}\par

    由 $x_M=x_P=(x_1+x_2)/2$，取抛物线上同横坐标的点
    \[
        N\left(\frac{x_1+x_2}{2},\frac{(x_1+x_2)^2}{8p}\right).
    \]
    $P,M,N$ 的横坐标相同。再计算
    \[
    \begin{aligned}
        y_N-y_P
        &=\frac{(x_1+x_2)^2-4x_1x_2}{8p}
          =\frac{(x_1-x_2)^2}{8p},\\
        y_M-y_N
        &=\frac{2(x_1^2+x_2^2)-(x_1+x_2)^2}{8p}
          =\frac{(x_1-x_2)^2}{8p}.
    \end{aligned}
    \]
    两者相等且为正，故点的顺序为 $P,N,M$，并且
    \[
        \boxed{PN=NM=\frac{PM}{2}}.
    \]

    \medskip
    \textbf{\textcolor{CProof}{四法比较（用于教师讲解与学生选择）}}
    \begin{itemize}[leftmargin=*, itemsep=0.4em, topsep=0.35em]
        \item \textbf{门槛最低：}切线点式 + 中点 + 两个普通三角形面积；最适合首次理解。
        \item \textbf{速度最快：}切线相减、平方差、中线面积三步；最适合考场直接代数计算。
        \item \textbf{步骤最少：}令横坐标为 $x_0\pm t$，两大步完成；适合总结对称结构。
        \item \textbf{常规通用：}从相切条件到“弦长$\times$点线距离”；适合迁移到一般直线与曲线的面积题。
    \end{itemize}
\end{proofbox}
